\section{A 0-Hecke action on chains of the composition poset}

\label{sec:0-hecke_operation_and_chains}

In~\cref{thm:sct_and_chains} a bijection between saturated chains in the composition poset $\mc L_c$ and standard composition tableaux is given. 
In this section we study the 0-Hecke action on these chains induced by this bijection. The main result of this section,~\cref{thm:support_of_column_word_and_shape}, will be useful in~\cref{sec:endomorphism_ring_of_0-Hecke_modules}. We begin with some notation.

\begin{defi}
	\label{thm:parts_of_chains_in_L_c}
	Let~$T$ be an SCT of shape~$\ab$ and size~$n$, $m\in [0,n]$ and ${\beta = \alpha^n \lessdot_c \alpha^{n-1} \lessdot_c \cdots \lessdot_c \alpha^0 =\alpha}$ the chain in $\mathcal L_c$ corresponding to~$T$. 
	 The SCT of shape $\alpha^m \skc \beta$ corresponding to the chain $\alpha^n \lessdot_c \alpha^{n-1}\lessdot_c \cdots \lessdot_c \alpha^m$ is denoted by $T^{>m}$.
\end{defi}

\begin{exam}
	For 
	$T=\begin{ytableau}
	1 \\
	*(gray) & *(gray) & 3 \\
	*(gray) & 2
	\end{ytableau}$ 
	we have
	$ 
	T^{>2} = \begin{ytableau} 
	*(gray) & *(gray) & 1 \\
	*(gray) \end{ytableau}.
	$	
\end{exam}

The following Lemma shows how we can obtain $T^{>m}$ directly from~$T$.

\begin{lemma}
	\label{thm:cells_and_entries_of_restricted_tableaux}
	Let~$T$ be an SCT of size~$n$ and shape~$\ab$, $\beta = \alpha^n \lessdot_c \alpha^{n-1} \lessdot_c \cdots \lessdot_c \alpha^0 = \alpha$ the chain in $\mathcal{L}_c$ corresponding to~$T$ and $m\in [0,n]$.
	\begin{enumerate}
		\item\label{kemma3.3_1}	
		$\alpha^m = \osh(T^{>m})$.
		\item\label{kemma3.3_2} We obtain $T^{>m}$ from~$T$ by removing the cells containing $1,\dots, m$ and subtracting~$m$ from the remaining entries.
	\end{enumerate}
\end{lemma}

	\begin{proof}
		Part~\eqref{kemma3.3_1} is a immediate consequence of~\cref{thm:parts_of_chains_in_L_c}. 	
		By~\cref{thm:sct_and_chains}, we obtain $T^{>m}$ by successively adding cells with entries $n-m,\dots, 1$ to the inner shape $\beta$ at exactly the same positions where we would add $n,\dots, m+1$ to $\beta$ in order to obtain~$T$ from its corresponding chain. This implies part~\eqref{kemma3.3_2}.
	\end{proof}

With the first part of~\cref{thm:cells_and_entries_of_restricted_tableaux} we can access the compositions within a chain of a given SCT. We use the following preorder to describe how the 0-Hecke action affects these compositions.

\begin{defi}\ \\*[-1.2em]

	\begin{enumerate}
		\item For a composition $\alpha = \parts{\alpha}{l}$ of $n$ and $j\in \N$ we define $|\alpha|_j = \# \left\{ i \in [l] \mid \alpha_i\geq j \right\}$. 
		\item
		 On the set of compositions of size~$n$ we define the preorder $\dleq$ by
		\begin{align*}
			\alpha \dleq \beta \iff \sum_{j=1}^k |\beta|_j \leq \sum_{j=1}^k |\alpha|_j \text{ for all $k \geq 1$.}
		\end{align*} 
		 Moreover, set $\alpha \dlneq \beta \iff \alpha \dleq \beta$ and $\alpha \neq \beta$.
	\end{enumerate}
\end{defi}

Note that $|\alpha|_j$ is the number of cells in the $j$th column of the diagram of $\alpha$.
Obviously $\dleq$ is reflexive and transitive. It is not antisymmetric since for example $(2,1) \dleq (1,2)$ and $(1,2) \dleq (2,1)$. In general, for $\alpha, \beta\vDash n$ we have
\begin{align*}
\alpha \dleq \beta \text{ and } 
\beta \dleq \alpha \iff |\alpha|_j = |\beta|_j \text{ for all $j=1,2,\dots$} \iff \widetilde \alpha = \widetilde \beta.
\end{align*}	

\begin{exam}
	\[
	\ydiagram{2,2} \dlneq \ydiagram{3,1}
	\qquad \qquad 
	\ydiagram{1,2,2}
	\dlneq
	\begin{ytableau}
	\hfil &&\\
	&
	\end{ytableau}		
	\]	
\end{exam}


	If we restrict $\dleq$ to partitions, we obtain the well known dominance order appearing, for example, in~\cite{Stanley1999}. However, $\dleq$ on partitions may seem to be reversed to the dominance order. This is because in the definition above we are considering the number of cells in columns rather than in rows as usual.







\begin{lemma}
		\label{thm:dominance_and_hecke_covering}
	Let~$\ab$ be a skew composition of size~$n$ and $T_1,T_2\in \SCT(\ab)$ be such that $T_2 = \pi_i T_1$ for an $i\in \nAD(T_1)$.
	 Then
	\begin{alignat*}{32}
		\osh(T_2^{>i}) &\dlneq \osh (T_1^{>i}), \\
		\osh(T_2^{>m}) &= \osh (T_1^{>m}) \text{ for } m \in [0,n], m\neq i.
	\end{alignat*}
\end{lemma}

\begin{proof}
	Assume $T_1,T_2$ and~$i$ as in the assertion. Then~$T_2$ is the tableau obtained from~$T_1$ by swapping the entries~$i$ and~$i+1$ of~$T_1$. Let $m\in[0,n]$.
	 
	If $m\neq i$ then either $\set{i,i+1} \subseteq [1,m]$ or $\set{i,i+1} \cap [1,m]= \emptyset$.
	Therefore, $T_1^{-1}([1,m]) = T_2^{-1} ([1,m])$ and so from the perspective of~\cref{thm:cells_and_entries_of_restricted_tableaux} we remove the same set of cells from~$T_1$ to obtain $T_1^{>m}$ as we remove from~$T_2$ to obtain $T_2^{>m}$. That is,	$\sh(T^{>m}_1) = \sh(T^{>m}_2)$.
	
	If $m = i$, set $(r_k,c_k) \coloneqq T_1^{-1}(k)$ for $k= i,i+1$, $\gamma_1 \coloneqq \osh(T^{>i}_1)$ and $\gamma_2 \coloneqq \osh(T^{>i}_2)$. 
	We assume that all composition diagrams appearing here are moved to the bottom of $\alpha$.
	Observe that as $T_2 = \fs_i T_1$, one obtains $\sh(T^{>i}_2)$ from $\sh(T^{>i}_1)$ by moving the cell $(r_{i+1}, c_{i+1})$ to the position $(r_i,c_i)$. Since $\ish(T^{>i}_2) = \beta = \ish(T^{>i}_1)$, we obtain $\gamma_2$ from $\gamma_1$ by this movement.
	 Moreover, $i\in \nAD(T_1)$ implies $c_i < c_{i+1}$. That is, we obtain $\gamma_2$ from $\gamma_1$ by moving a cell strictly to the left. By the definition of $\dleq$, this means that $\gamma_2 \dlneq \gamma_1$.
\end{proof}

 \begin{exam} 
 	The~$\He$-action on tableaux and the corresponding chains of the composition poset is shown below.
 	\begin{align*}
 	\begin{array}{c|ccccccc}
 	T & \osh(T^{>3}) &-& \osh(T^{>2}) &-& \osh(T^{>1}) &-& \osh(T^{>0})\\ \hline
 	\begin{ytableau}
 	1 \\
 	*(gray) & *(gray) & 3 \\
 	*(gray) & 2
 	\end{ytableau}
 	& \ydiagram{2,1} 
 	&-& \ydiagram{3,1} 
 	&-& \ydiagram{3,2} 
 	&-& \ydiagram{1,3,2} \\
 	\downarrow \pi_2 &&& \downarrow \pi_2 
 	\\
 	\begin{ytableau}
 	1 \\
 	*(gray) & *(gray) & 2 \\
 	*(gray) & 3
 	\end{ytableau}
 	& \ydiagram{2,1} 
 	&-& \ydiagram{2,2} 
 	&-& \ydiagram{3,2} 
 	&-& \ydiagram{1,3,2} \\
 	\downarrow \pi_1 &&&&& \downarrow \pi_1
 	\\
 	\begin{ytableau}
 	2 \\
 	*(gray) & *(gray) & 1 \\
 	*(gray) & 3
 	\end{ytableau}
 	& \ydiagram{2,1} 
 	&-& \ydiagram{2,2} 
 	&-& \ydiagram{1,2,2} 
 	&-& \ydiagram{1,3,2} \\	
 	\end{array}
 	\end{align*}
 \end{exam}
Let~$\ab$ be a skew composition, $E\in \mc E(\ab)$ and $T_1,T_2 \in E$ be such that $T_1\pleq T_2$.
Recall that for each saturated chain from~$T_1$ to~$T_2$ in~$E$ the index set of the 0-Hecke operators establishing the covering relations within the chain is $\supp(\col_{T_2}\col^{-1}_{T_1})$.
As a consequence of~\cref{thm:dominance_and_hecke_covering} we obtain a criterion for determining whether an operator $\pi_i$ appears in the saturated chains from~$T_1$ to~$T_2$ or not.

\begin{prop}
	\label{thm:support_of_column_word_and_shape}
	Let~$\ab$ be a skew composition of size~$n$, $i\in [n-1], E\in \mc E(\ab)$ and $T_1,T_2\in E$ be such that $T_1\pleq T_2$. 
	Then
	$i\in \supp(\col_{T_2}\col^{-1}_{T_1})$
	if and only if
	$\sh(T_2^{>i}) \neq \sh(T_1^{>i})$.
\end{prop}

\begin{proof}	 
	\cref{thm:dominance_and_hecke_covering} applied to each covering relation in a saturated chain from~$T_1$ to~$T_2$ in~$E$ and the fact that $\dleq$ is a preorder imply 
	\begin{align*}
		i\in \supp(\col_{T_2}\col^{-1}_{T_1}) \iff \osh(T_2^{>i}) \neq \osh(T_1^{>i}).
	\end{align*}
	From this we obtain the claim since $\ish(T^{>i}_1) = \beta = \ish(T^{>i}_2)$.
\end{proof}

\section{The endomorphism ring of \texorpdfstring{$\bs S_{\alpha,E}$}{S\_a,E}}
\label{sec:endomorphism_ring_of_0-Hecke_modules}

 For each $\alpha \vDash n$ there is an equivalence class $E_\alpha\in \mc E(\alpha)$ such that for all $T \in E_\alpha$ the entries increase in each column from top to bottom~\cite[Section~8]{Tewari2015}. 
In~\cite{Tewari2015}, Tewari and van Willigenburg showed that $\bs S_{\alpha,E_\alpha}$ is indecomposable.

In this section, we show for all $E\in \mc E(\alpha)$ that $\End_{\He}(\se) = \K\id$ and hence~$\se$ is indecomposable; this extends the result of Tewari and van Willigenburg to the general case.
By~\cref{thm:decomposition_in_equivalence_classes} we then have the desired decomposition of~$\sa$.
In contrast, skew modules $\sse$ can be decomposable (see~\cref{exa:decomposable_module} at the end of the section).

We fix some notation that we use in the entire section unless stated otherwise.
Let $\alpha\vDash n$, $E\in \mc E(\alpha)$ and $T_0 \coloneqq T_{0,E}$ be the source tableau of~$E$.
Moreover, let $f\in \End_{\He}(\se)$, $v\coloneqq f(T_0)$ and $v = \sum_{T\in E} a_T T$ be the expansion of~$v$ in the $\K$-basis~$E$.
Since~$\se$ is cyclically generated by~$T_0$,~$f$ is already determined by~$v$.
The \emph{support} of~$v$ is given by $\supp(v)=\set{T\in E \mid a_T \neq 0}$.
Our goal is to show that~$T_0$ is the only tableau that may occur in $\supp(v)$ since then $f = a_{T_0} \id \in \K\id$.
We begin with a property holding for $\supp(v)$ that also appeared in the proof of~\cite[Theorem~7.8]{Tewari2015}.

\begin{lemma}
	\label{thm:descent_set_of_support_of_v}
	If $T\in \supp(v)$ then $\D(T) \subseteq \D(T_0)$.
\end{lemma}
\begin{proof} Let $T_*\in E$ be such that $\D(T_*)\not \subseteq \D(T_0)$. 
	Then there is an $i\in \D(T_*)\cap \DC(T_0)$. 
	 Because $i\in \DC(T_0)$, $\pi_i v = f(\pi_i T_0) = v$. Thus, $a_{T_*}$ is the coefficient of~$T_*$ in $\pi_i v = \sum_{T \in E} a_T \pi_i T$. But this coefficient is~$0$ since $\pi_iT \neq T_*$ for all $T\in E$. To see this, assume that there is a $T\in E$ such that $\pi_i T = T_*$. Then we obtain a contradiction as
	 \[
	 T_* \neq \pi_i T_* = \pi_i^2 T = \pi_i T = T_*. \qedhere
	 \]	
\end{proof}

Thanks to~\cref{thm:descent_set_of_support_of_v} it remains to show $a_T = 0$ for all $T\in E$ such that $T\neq T_0$ and $D(T)\subseteq D(T_0)$. So fix such a tableau~$T$.
In order to determine~$a_T$, we will use a 0-Hecke operator $\pi_\sigma$ where $\sigma~=~s_{j-1}\cdots s_i$ and~$i$ and~$j$ are given by
\begin{equation}
	 \label{eq:def_of_i_and_j}
	 \begin{aligned}
	 i &= \max \big\{ k \in [n] \mid T^{-1}(k) \neq T^{-1}_0(k) \big\}, \\
	 j &= \min \big\{k \in [n] \mid k>i \text{ and } i\att_{T_0} k \big\}.
	 \end{aligned}
\end{equation}
That is,~$i$ is the greatest entry whose position in~$T$ differs from that in~$T_0$ and~$j$ is the smallest entry in~$T_0$ which is greater than~$i$ and attacked by~$i$ in~$T_0$. 
At this point it is not clear that~$j$ is well defined since the defining set could be empty.
However, the next two lemmas will show that there always exists an element in this set.

\begin{samepage}
\begin{exam}
	Consider the equivalence class~$E$ from~\cref{fig:example_for_straight_E}. Then $T_0 = T_{0,E}$ and there is exactly one other tableau~$T$ in~$E$ with $\D(T)\subseteq \D(T_0)$:
	\begin{align*}
	T_0
	= 
	\begin{ytableau} 
	1 \\ 
	6 & 5 & \bs 4 & 3 \\ 
	8 & 7 & \bs 2 \\ 
	\end{ytableau} 
	\overset{\pi_1}{\longrightarrow}
	T = 
	\begin{ytableau} 
	\bs 2 \\ 
	6 & 5 & \bs 4 & 3 \\ 
	8 & 7 & 1 \\ 
	\end{ytableau}
	\end{align*}
	Defining~$i$ and~$j$ for~$T$ as in~\eqref{eq:def_of_i_and_j}, we obtain $i=2$ and $j=4$.
	Note that $2\in D(T_0)$. This property holds in general by the following result.
\end{exam}	
\end{samepage}

\begin{lemma} 	
	\label{thm:i_is_descent}
	Let $T\in E$ be such that $T\neq T_0$ and $\D(T)\subseteq \D(T_0)$ and set
	\begin{align*}
	i = \max \big\{k \in [n] \mid T^{-1}(k) \neq T^{-1}_0(k)\big\}.
	\end{align*}
	Then $i\in \D(T_0)$.
\end{lemma}

\begin{proof}
	Let~$T$,~$T_0$ and~$i$ be given as in the assertion. 
	We introduce indices such that
	$D(T_0)= \set{ d_1 < d_2 < \cdots < d_m}$ and set $d_0 \coloneqq 0, d_{m+1} \coloneqq n$. Moreover,
	 define $I_k\coloneqq [d_{k-1}+1 , d_k]$ for $k= 1, \ldots, m+1$. 
	 Recall that since~$T_0$ is a source tableau, $\DC(T_0)= \ND(T_0)$ by~\cref{thm:source_and_sink_tableau}. 
	 That is, $a+1$ is the left neighbor of $a$ for each ascent~$a$ of~$T_0$. Therefore, we have $I_k\setminus \set{d_k} \subset \ND(T_0)$ and conclude that $T^{-1}_0(I_k)$ is a connected horizontal strip (a one-row diagram which contains all cells between its leftmost and rightmost cell) for $k = 1,\dots,m+1$.
	 
	Set $\cell_k \coloneqq T_0^{-1}(k)$ for $k=1,\dots, n$ and let~$x$ be the index such that $T(\cell_x)=~i$. 
	Since~$T_0$ and~$T$ are straight, the ordering conditions of standard composition tableaux imply $T^{-1}(n)= (\ell(\alpha), 1) = T_0^{-1}(n)$. Therefore $i\neq n$ and we now show $i\notin \DC(T_0)$.	
	
	Assume for sake of contradiction that $i\in \DC(T_0)$.
	Let $l\in [m+1]$ be such that $i\in I_l$. 
	Since $i \in \DC(T_0)$, $i<d_l$ and $i+1\in I_l$. 
	The strip $T^{-1}_0(I_l)$ looks as follows:
	\begin{align}
	\label{eq:horizontal_strip_of_i}
	\cell_{d_{l}} \cell_{d_{l} -1} \cdots \cell_{i+1} \cell_i \cdots \cell_{d_{l-1} + 1}.
	\end{align}	

	By choice of~$i$, we have 
	\begin{align}
	\label{eq:entries_of_T_greater_i}
	\text{$T(\cell_k) = k$ for $k=i+1,\dots, n$ and 	$T(\cell_i) < i$.}
	\end{align}
	Since entries decrease in rows of~$T$,~\eqref{eq:horizontal_strip_of_i} implies 
	\begin{align}
	\label{eq:entries_of_T_in_horizontal_strip_smaller_i}
	T(\cell_k) < i \text{ for } k=d_{l-1}+1,\dots , i. 
	\end{align}
	Combining~\eqref{eq:entries_of_T_greater_i} and~\eqref{eq:entries_of_T_in_horizontal_strip_smaller_i} we obtain
	\begin{align}
	\label{eq:bound_for_x}
	 x\leq d_{l-1}.
	\end{align}
	We deal with two cases depending on $c_T(i)$. 
	In both cases we will end up with a contradiction.
	
\begin{enonce*}[remark]{Case 1} 
$c_T(i) \leq c_{T_0}(d_{l-1}+1)$. It follows from $\D(T) \subseteq \D(T_0)$ that $i\in \DC(T)$ and thus $c_T(i+1) < c_T(i) $. 
	Using $c_{T_0}(i) = c_{T_0}(i+1) + 1 = c_{T}(i+1) + 1$, we obtain
	$c_{T_0}(i) \leq c_T(i) \leq c_{T_0}(d_{l-1}+1).$
	Then there is a $y\in [d_{l-1}+1, i]$ such that $\cell_x$ and $\cell_y$ are in the same column. 
	On the one hand, we obtain from~\eqref{eq:entries_of_T_in_horizontal_strip_smaller_i} that $T(\cell_y) < i = T(\cell_x)$. 
	On the other hand, the choice of $y$ and~\eqref{eq:bound_for_x} imply $y>d_{l-1} \geq x$ and hence $T_0(\cell_y) =y> x = T_0(\cell_x)$.	
	That is, in the column of $\cell_x$ and $\cell_y$ the relative order of entries in $T$ differs from that in~$T_0$.
	So $T\not \sim T_0$ which contradicts the assumption $T,T_0 \in E$.
	\end{enonce*}
\begin{enonce*}[remark]{Case 2} 
	$c_T(i) > c_{T_0}(d_{l-1}+1)$. This case is illustrated in~\cref{fig:i_is_descent}. Since by~\eqref{eq:bound_for_x} $x\leq d_{l-1}$, there is a $1\leq p\leq l-1$ such that $x\in I_p$.
	The leftmost cell of the connected horizontal strip $T^{-1}_0(I_p)$ is $\cell_{d_p}$. As entries decrease in rows of $T$ from left to right, we have $T(\cell_{d_p})\geq T(\cell_x) = i$. In addition, the choice of $p$ and~\eqref{eq:entries_of_T_greater_i} imply that $T(\cell_{d_p})\leq i$. Thus, $d_p=x$.
	
	From $d_p=x$ we obtain $d_p \neq d_{l-1}$ since
	\begin{align*}
	c_{T_0}(d_{l-1}) \leq c_{T_0}(d_{l-1}+1) < c_T(i) = c_{T_0}(d_p)
	\end{align*}
	where we use $d_{l-1}\in \D(T_0)$ for the first inequality.
	
	We claim that there exists an index $y\in [d_p +1, d_{l-1} -1]$ such that $\cell_y$ and $\cell_{d_p}$ are located in the same column.
	To prove the claim, assume for sake of contradiction that this is not the case. Then $d_p \in \D(T_0)$ implies $ c_{T_0}(d_p)< c_{T_0}(d_p+1)$. Thus, it follows from $\DC(T_0)= \ND(T_0)$ and induction that $c_{T_0}(d_p)< c_{T_0}(z)$ for all ${z\in[d_p+1,d_{l-1} -1]}$. As a consequence,
		\begin{align*}
		c_{T_0}(d_{l-1}) < c_{T_0}(d_p) < c_{T_0}(d_{l-1}-1).
		\end{align*}
		In other words, $d_{l-1}-1$ is an ascent of~$T_0$
		but $d_{l-1}$ is not the left neighbor of $d_{l-1}-1$.	
		 This is a contradiction to the fact that~$T_0$ is a source tableau and finishes the proof of the claim.
	
	Now, let $y$ be as claimed above. Then $y\in [d_p +1, d_{l-1} -1]$ and in particular $y\neq d_p= x$. Hence,~\eqref{eq:entries_of_T_greater_i} implies $T(\cell_y) <i$ and so $T(\cell_y) < i = T(\cell_{d_p})$ . 
 On the other hand, $y\in [d_p +1, d_{l-1} -1]$ yields $T_0(\cell_y) =y > d_p = T_0(\cell_{d_p})$. As in Case 1, this is a contradiction to $T,T_0\in E$.\qedhere
\end{enonce*}
\end{proof}

\begin{figure}[htb]\centering
	\begin{tikzpicture}[scale=.75, transform shape]
	
	\def \x{-3}
	\def\y{0}
	\node at (\x,\y) {$d_l$};
	\draw (\x-0.5,\y+0.5) rectangle (\x+0.5,\y-0.5);
	\node[below right] at (\x+0.5,\y-0.5) {\scriptsize $d_l$};
	
	\def \x{-1.5}
	\def\y{0}
	\node at (\x,\y) {$\dots$};
	
	\def \x{0}
	\def\y{0}
	\node at (\x,\y) {$i+1$};
	\draw (\x-0.5,\y+0.5) rectangle (\x+0.5,\y-0.5);
	\node[below right] at (\x+0.5,\y-0.5) {\scriptsize $i+1$};
	
	\def \x{1}
	\def\y{0}
	\node at (\x,\y) {$<i$};
	\draw (\x-0.5,\y+0.5) rectangle (\x+0.5,\y-0.5);
	\node[below right] at (\x+0.5,\y-0.5) {\scriptsize $i$};
	
	\def \x{3}
	\def\y{0}
	\node at (\x,\y) {$\dots$};
	
	\def \x{3}
	\def\y{2}
	\node at (\x,\y) {$<i$};
	\draw (\x-0.5,\y+0.5) rectangle (\x+0.5,\y-0.5);
	\node[below right] at (\x+0.5,\y-0.5) {\scriptsize $d_{l-1}$};
	
	\def \x{5}
	\def\y{0}
	\node at (\x,\y) {$<i$};
	\draw (\x-0.5,\y+0.5) rectangle (\x+0.5,\y-0.5);
	\node[below right] at (\x+0.5,\y-0.5) {\scriptsize $d_{l-1}+1$};
	
	\draw[loosely dashed] (\x+0.5,\y-3) -- (\x+0.5,\y+4.5);
	
	\def \x{7.5}
	\def\y{-2}
	\node at (\x,\y) {$i$};
	\draw (\x-0.5,\y+0.5) rectangle (\x+0.5,\y-0.5);
	\node[below right] at (\x+0.5,\y-0.5) {\scriptsize $d_{p}$};
	
	\def\y{3.5}
	\node at (\x,\y) {$<i$};
	\draw (\x-0.5,\y+0.5) rectangle (\x+0.5,\y-0.5);
	\node[below right] at (\x+0.5,\y-0.5) {\scriptsize $y$};
	
	\end{tikzpicture} 
	\caption{
		\label{fig:i_is_descent}
		An example for the positions of cells and entries in the tableau $T$ from Case 2 of the proof of~\cref{thm:i_is_descent}.}
\end{figure}

Note that the $i$ appearing in the following Lemma is not the same as in~\eqref{eq:def_of_i_and_j}.

\begin{lemma}
	\label{thm:attacked_entry_exists}
	For all $i\in \D(T_0)$ there exists $k\in T_0$ such that $k>i$ and $i \att_{T_0} k$.
\end{lemma}
\begin{proof}
	Let $i\in \D(T_0)$.
	Then $c_{T_0}(i) \leq c_{T_0}(i+1)$ and thus $r_{T_0}(i) \neq r_{T_0}(i+1)$.
	Since~$T_0$ is straight by assumption, the cell $(r_{T_0}(i+1), c_{T_0}(i))$ is contained in the shape of~$T_0$.
	Let $k$ be the entry of~$T_0$ in that cell.
	Then $i\att_{T_0} k$ and $k\geq i+1 $ because entries decrease in rows.
\end{proof}

Let $T$, $i$ and~$j$ as in~\eqref{eq:def_of_i_and_j}.~\cref{thm:i_is_descent} and~\cref{thm:attacked_entry_exists} now show that~$j$ is well defined. We proceed by considering the relative positions of~$i$ and $[i+1,j]$ first in~$T_0$ and then in~$T$. This will allow us to deduce useful properties of the operator $\pi_\sigma$ to be defined in~\cref{thm:end:annihilator_for_T_0}. In the following Lemma,~$i$ is slightly more general than in~\eqref{eq:def_of_i_and_j}.

\begin{lemma}
	\label{thm:position_of_entries_smaller_as_attacked_entry_in_T_0}
	Let $i\in \D(T_0)$ and set $j = \min \{k \in [n] \mid k>i \text{ and } i\att_{T_0} k \}$.
	Then~$j$ is well defined and in~$T_0$~$i$ is located strictly left of $[i+1,j-1]$ and does not attack $[i+1,j-1]$.
\end{lemma}

We illustrate~\cref{thm:position_of_entries_smaller_as_attacked_entry_in_T_0} before we prove it.

	\begin{exam} For the source tableau from above
		\[
			T_0
			= 
\begin{ytableau} 
1 \\ 
6 & 5 & \bs4 & \bs3 \\ 
8 & 7 & \bs2 \\ 
\end{ytableau} 
		\]
		and $i=2\in \D(T_0)$ we have $j = 4 =\min \{k \in [n] \mid k>i \text{ and } i\att_{T_0} k \}$ and $\set{3} = [i+1,j-1]$. Note $2\att_{T_0} 4$ but $2\not \att_{T_0} 3$.
	\end{exam}

	\begin{proof}[Proof of~\cref{thm:position_of_entries_smaller_as_attacked_entry_in_T_0}]
		First, it follows from~\cref{thm:attacked_entry_exists} that~$j$ is well defined. 
		We set $I =: [i+1,j-1]$ and $c_l \coloneqq c_{T_0}(l)$ for $l \in T_0$.
		By the minimality of~$j$, we have $i\not \att_{T_0} I$.
		It remains to show that~$i$ is strictly left of~$I$ or equivalently that $c_i < c_l$ for all $l\in I$.
		We may assume $I\neq \emptyset$ and use an induction argument to show this.
		
		We begin with~$i+1$, the minimum of~$I$.
		 Since $i\in \D(T_0)$, $c_i \leq c_{i+1}$. Moreover, $i+1\in I$ implies $i\not \att_{T_0} i+1$ and consequently $c_i < c_{i+1}$.
		
		Now, let $l\in I$ such that $l>i+1$ and $c_i < c_{l-1}$.
		If $l-1\in \D(T_0)$ then $c_i < c_{l-1} \leq c_l$.
		If $l-1\in \DC(T_0)$ then $l-1\in \ND(T_0)$ as~$T_0$ is a source tableau.
		Thus $c_l= c_{l-1}-1$ and $c_i \leq c_l$. Furthermore $c_i \neq c_l$ since $i\not \att_{T_0} I \ni l$.
		Hence, $c_i< c_l$.
			 \qedhere
	\end{proof}

Let~$T,i$ and~$j$ as in~\eqref{eq:def_of_i_and_j}.
By definition, $i$ attacks~$j$ in~$T_0$.
In contrast, the next Lemma shows that~$i$ does not attack~$j$ in~$T$.
There,~$i$ and~$j$ are defined as in~\eqref{eq:def_of_i_and_j}.

\begin{lemma}
	\label{thm:position_of_entries_smaller_as_attacked_entry_in_T}
	Let $T\in E$ be such that $T\neq T_0$ and $\D(T) \subseteq \D(T_0)$. Define
	\begin{align*}
		i &= \max \big\{k \in [n] \mid T^{-1}(k) \neq T^{-1}_0(k)\big\}, \\
		j &= \min \big\{k \in [n] \mid k>i \text{ and } i\att_{T_0} k\big\}.
	\end{align*}
	Then~$i$ and~$j$ are well defined and in~$T$ $i$ appears strictly left of $[i+1,j]$ and does not attack $[i+1,j]$. 
\end{lemma}

We first give an example and then the proof of~\cref{thm:position_of_entries_smaller_as_attacked_entry_in_T}\!.

\begin{samepage}
\begin{exam} Recall that in our running example $i=2$ and $j=4$ when defined for
	\begin{align*}
		T &= 
		\begin{ytableau} 
			\bs 2 \\ 
			6 & 5 & \bs4 & \bs3 \\ 
			8 & 7 & 1 \\ 
		\end{ytableau}
	\end{align*} 
	 as in~\cref{thm:position_of_entries_smaller_as_attacked_entry_in_T}. Then
 $[i+1,j] = \set{3,4}$ and $2\not \att_T \set{3,4}$.
\end{exam}
\end{samepage}


\begin{proof}[Proof of~\cref{thm:position_of_entries_smaller_as_attacked_entry_in_T}] 
	\cref{thm:i_is_descent} yields $i\in \D(T_0)$ and so~\cref{thm:attacked_entry_exists} ensures that~$j$ is well defined. Set $\sigma \coloneqq \col_T\col^{-1}_{T_0}$, $\cell_k \coloneqq T_0^{-1}(k)$ for $k=1,\dots, n$ and let~$x$ be the index such that $T(\cell_x)= i$. 

	 By choice of~$i$, we have $T^{>i} = T_0^{>i}$. So, $\sh (T^{>k}) = \sh (T_0^{>k})$ for $k = i ,\dots, n$. Hence, from~\cref{thm:support_of_column_word_and_shape} we obtain 
		\begin{align}
		\label{eq:support_of_path_to_T}
			\supp(\sigma) \subseteq [i-1].	
		\end{align}
	Let $\fs_{i_p} \cdots \fs_{i_1}$ be a reduced word for $\sigma$.
	Then $T = \pi_{i_p} \cdots \pi_{i_1} T_0$. From~\eqref{eq:support_of_path_to_T} we have $i_q \neq i$ for $q=1,\dots, p$.
	Moreover, at least one $\pi_{i_q}$ has to move~$i$ because the position of~$i$ in~$T$ differs from its position in~$T_0$. Hence, there is a~$q$ such that $i_q = i-1$ since $\pi_{i-1}$ and $\pi_i$ are the only operators that are able to move~$i$.
	For two standard composition tableaux~$T_1$ and~$T_2$ such that $T_2 = \pi_{i-1}T_1= \fs_{i-1} T_1$ we have that $i-1\in \nAD (T_1)$ and thus $T_2^{-1}(i)$ is left of $T_1^{-1}(i)$ and $T_2^{-1}(i) \not \att T_1^{-1}(i)$. So, by applying $\pi_{i_p} \cdots \pi_{i_1}$ to~$T_0$,~$i$ is moved (possibly multiple times) strictly to the left into a cell that does not attack $\cell_i$. That is,
	\begin{align}
		\label{eq:postitoin_of_i_in_T_and_T_0}
		\text{$\cell_x$ is located strictly left of $\cell_i$ and $\cell_x\not \att \cell_i$.}
	\end{align}
	It follows from the choice of~$i$ that the elements of $[i+1,j-1]$ have the same position in~$T$ and~$T_0$. By combining~\eqref{eq:postitoin_of_i_in_T_and_T_0} and~\cref{thm:position_of_entries_smaller_as_attacked_entry_in_T_0} we obtain:
	\begin{align}
	\text{In~$T$~$i$ is located strictly left of $[i+1,j-1]$ and $i\not \att_T [i+1,j-1]$.}
	\end{align}
	Recall that~$j$ has the same position in~$T$ and~$T_0$. 
	It follows from~\eqref{eq:postitoin_of_i_in_T_and_T_0} and $i\att_{T_0} j$ that $c_T(i) < c_{T_0}(i)\leq c_{T_0}(j)$. Thus,~$i$ is strictly left of~$j$ in~$T$.
	
	It remains to show $i\not \att_T j$.
	 We have either $c_{T_0}(j)=c_{T_0}(i)+1$ or $c_{T_0}(j)=c_{T_0}(i)$ since $i\att_{T_0} j$.
	
\begin{enonce*}[remark]{Case 1} $c_{T_0}(j)=c_{T_0}(i)+1$. Then~\eqref{eq:postitoin_of_i_in_T_and_T_0} implies $c_T(i) < c_{T_0}(i)< c_{T_0}(j) = c_T(j)$ and so $i\not \att_T j$. 
\end{enonce*}

\begin{enonce*}[remark]{Case 2} $c_{T_0}(j)=c_{T_0}(i)$. If $c_T(i) < c_{T_0}(i)-1$ then $c_T(i) < c_T(j)-1$ and so $i\natt_T j$. If $c_T(i) = c_{T_0}(i)-1$ then~$i$ and~$j$ appear in adjacent columns of~$T$ and for $i\natt_T j$ we have to show that $r_T(j) < r_T(i)$.
	 On the one hand, we have $1\leq c_T(i) <c_{T_0}(i)$ so that~$i$ has a left neighbor $t>i$ in~$T_0$.
	 In addition, from the first statement of~\cref{thm:position_of_entries_smaller_as_attacked_entry_in_T_0} and $c_{T_0}(j)=c_{T_0}(i)$ we obtain that~$i$ is weakly left of $[i+1,j]$ in~$T_0$.
	 Thus, $t>j$ and hence $r_{T_0}(j) < r_{T_0}(i)$ because otherwise~$t,i$ and~$j$ would violate the triple rule in~$T_0$.
	 On the other hand, $c_T(i) = c_{T_0}(i)-1$ and $i\not \att_T \cell_i$ imply $r_{T_0}(i) < r_T(i)$. All in all, $r_{T}(j) = r_{T_0}(j) < r_{T_0}(i) < r_T(i)$ and thus $i \not \att_T j$.\qedhere
\end{enonce*}
\end{proof}

Next, we prove useful properties of the operators $\pi_\sigma$ mentioned already above in~\eqref{eq:def_of_i_and_j}.

\begin{lemma}
	\label{thm:end:annihilator_for_T_0}
	Keep the notation of~\cref{thm:position_of_entries_smaller_as_attacked_entry_in_T} and set $\sigma = \fs_{j-1} \cdots \fs_{i+1} \fs_i$. Then
	\begin{enumerate}
		\item\label{lemma4.9_1} $\pi_\sigma T_0 = 0$,
		\item\label{lemma4.9_2} $\pi_\sigma T \in E$,
		\item\label{lemma4.9_3} $\sigma = \col_{\pi_\sigma T} \col^{-1}_{T}$.
	\end{enumerate}
\end{lemma}



\begin{proof}
	First observe that $\fs_{j-1} \cdots \fs_{i+1} \fs_i$ is a reduced word, \ie $\pi_\sigma = \pi_{j-1} \cdots \pi_{i+1}\pi_i$. Set $\cell_k = T_0^{-1}(k)$ for $k=1,\dots, n$.
	
	We consider~$T_0$. Set $T' = \pi_{j-2} \cdots \pi_{i+1}\pi_i T_0$. 
	We can apply~\cref{thm:mutliple_flips_same_cell} in~$T_0$ to~$i$ and $[i+1,j-1]$ because of~\cref{thm:i_is_descent} and~\cref{thm:position_of_entries_smaller_as_attacked_entry_in_T_0}. By doing this, we obtain that $T'\in E$ and $T'(\cell_i) = j-1$. In addition, $T'(\cell_j) = T_0(\cell_j) = j$ as none of the operators $\pi_{j-2}, \ldots, \pi_{i+1}$ moves~$j$. Recall that~$j$ is defined such that $\cell_i \att \cell_j$. Thus $j-1\in \AD(T')$ and $\pi_\sigma T_0 = \pi_{j-1}T' = 0$.
	
	
	Now consider~$T$. Because of~\cref{thm:position_of_entries_smaller_as_attacked_entry_in_T}, we can apply~\cref{thm:mutliple_flips_same_cell} in~$T$ on~$i$ and $[i+1,j]$. This immediately gives us~\eqref{lemma4.9_2} and~\eqref{lemma4.9_3}.
\end{proof}

\begin{exam} Continuing our running example, we have $i = 2$, $j = 4$ and $\pi_\sigma=\pi_3\pi_2$. Moreover,
	\begin{align*}
		T_0 &=
		\begin{ytableau} 
			1 \\ 
			6 & 5 & 4 & 3 \\ 
			8 & 7 & 2 \\ 
		\end{ytableau} 	
		\overset{\pi_2}{\longrightarrow}
		\begin{ytableau} 
			1 \\ 
			6 & 5 & 4 & 2 \\ 
			8 & 7 & 3 \\ 
		\end{ytableau} 	
		\overset{\pi_3}{\longrightarrow}
		0,
		\\
		T &= 
		\begin{ytableau} 
			2 \\ 
			6 & 5 & 4 & 3 \\ 
			8 & 7 & 1 \\ 
		\end{ytableau} 
		\overset{\pi_2}{\longrightarrow}
		\begin{ytableau}
			3 \\ 
			6 & 5 & 4 & 2 \\ 
			8 & 7 & 1 \\ 
		\end{ytableau}
		\overset{\pi_3}{\longrightarrow}
		\begin{ytableau}
			4 \\ 
			6 & 5 & 3 & 2 \\ 
			8 & 7 & 1 \\ 
		\end{ytableau}.
	\end{align*}
\end{exam}

We are ready to prove the main result of this paper now.

